\chapter{Discussion} \label{chap:discussion}

\begin{quote}
    \textit{The view that machines cannot give rise to surprises is due, I believe, to a fallacy to which philosophers and mathematicians are particularly subject. This is the assumption that as soon as a fact is presented to a mind, all consequences of that fact spring into the mind simultaneously with it. It is a very useful assumption under many circumstances, but one too easily forgets that it is false. A natural consequence of doing so is that one then assumes that there is no virtue in the mere working out of consequences from data and general principles.}

    \begin{flushright}
        --- Alan Turing (1950) \nocite{turing1950computing}
    \end{flushright}
\end{quote}

\vspace{1em}

As this dissertation reaches its final chapter, we take a moment to look both backward and forward. Throughout this work, we investigate the application of generative models to large-scale inference problems in physics, focusing in particular on systems whose states evolve over time. Such dynamical systems are ubiquitous in nature, science, and engineering and thus constitute a key research area for technological and societal progress. Across the five preceding chapters, we explore multiple facets of probabilistic modeling and dynamical systems, including state estimation, forecasting, reduced-order modeling, and learning from corrupted observations. We now pause to reflect on what has been accomplished, what has been learned, and what lies ahead.

\vspace{-1ex}
\section{Impact}

The cornerstone of our work, introduced in Chapter \ref{chap:sda} and extended in Chapters \ref{chap:sda-2lqg} and \ref{chap:appa}, is a novel framework that combines diffusion models~\cite{sohl-dickstein2015deep, ho2020denoising} and training-free guidance methods~\cite{chung2023diffusion, song2023pseudoinverseguided, boys2023tweedie} to estimate the state of dynamical systems from a series of observations over time; a task known as data assimilation (DA)~\cite{ledimet1986variational, tremolet2006accounting, hamill2006ensemblebased, evensen2009data, carrassi2018data}. Unlike classical DA methods, our method does not require simulating or differentiating through a hand-crafted physical model at inference, and is able to assimilate any observation. While previous studies combine machine learning methods and DA~\cite{mack2020attentionbased, brajard2020combining, brajard2021combining, arcucci2021deep, fablet2021learning, frerix2021variational, yeung2022physicsinformed, andrychowicz2023deep, cheng2023machine}, we are the first to consider generative models, especially diffusion models, in this context. Following the publication of SDA~\cite{rozet2023scorebased}, many studies applied generative models to DA~\cite{huang2024diffda, foroumandi2025harnessing, martin2025generative, manshausen2025generative, schmidt2025generative, shen2025conditional, sun2025alignda, chakraborty2025multimodal, xu2025fuxida}. We are confident that this trend will continue as generative modeling offers a principled way to handle uncertainty, unlike artificially perturbing observations to obtain an ensemble of solutions~\cite{liu2013ensemblebased, hoteit2015mitigating, auligne2016ensemble, bannister2017review}.

While SDA establishes diffusion models as a promising alternative to classical DA methods, the data it relies on during training -- fully observed state sequences -- is rarely accessible. In practice, only incomplete and noisy observations of the states are available, especially in earth and space sciences where the systems of interest can only be probed superficially. This limitation motivates our second contribution, presented in Chapter \ref{chap:diem}, an adaptation of the expectation-maximization algorithm~\cite{hartley1958maximum, dempster1977maximum, wu1983convergence, mclachlan2007em, balakrishnan2017statistical} tailored to diffusion models. With our method, named diffusion-based expectation-maximization (DiEM), it becomes possible to train or tune diffusion models when only corrupted data is available. Remarkably, the learned diffusion model can later be used as prior in Bayesian inference problems, such as gravitational lensing inversion~\cite{warren2003semilinear, morningstar2019datadriven, mishra-sharma2022strong}, accelerated MRI~\cite{wang2016accelerating, hammernik2018learning, han2020kspace, zbontar2018fastmri, knoll2020fastmri}, or ... data assimilation~\cite{ledimet1986variational, tremolet2006accounting, hamill2006ensemblebased, evensen2009data, carrassi2018data}. Although this research question has been approached from different directions~\cite{aali2023solving, daras2023ambient, daras2024consistent, kawar2024gsurebased}, learning from corrupted data remains an underexplored research area, with virtually endless applications and unparalleled potential for impact.

Orthogonal to data availability, computational cost is another limitation of SDA, which becomes prohibitive for long trajectories and high-dimensional states, although it remains cheaper than 4D-Var~\cite{tremolet2006accounting, tremolet2007modelerror}. Our last contribution, unifying the findings of Chapters \ref{chap:lola} and \ref{chap:appa}, is to demonstrate that modeling dynamical systems in a compressed latent representation is a sensible approach to reduce training and inference cost. Latent models are not only more efficient, but can also improve the accuracy of emulation and assimilation compared to physical-space models. In the context of numerical weather prediction, this efficiency could enable assimilation of a larger volume of observations or include physical processes that are currently omitted due to their computational burden, such as the interaction between the atmosphere and oceans.

\vspace{-1ex}
\section{Limitations \& Perspectives}

While the results presented in this dissertation are promising, several methodological and experimental limitations remain. Each of these limitations represents an exciting opportunity for future research, and addressing them will be essential to accelerate the adoption of generative modeling and, more broadly, deep learning by domain scientists. We organize our discussion around three axes: long-term dependencies, posterior validation, and learning from observations.

\vspace{-1ex}
\subsection*{The frontier of long-term assimilation}

A first limitation concerns SDA's ability to capture long-term dependencies. The pseudo-blanket approximation \eqref{eq:sda-pseudo-blanket}, although accurate for high signal-to-noise regimes ($\sigma_t \to 0$), introduces errors as the noise magnitude increases. In retrospect, evaluating this approximation on simplified problems would have provided valuable insight. Later studies~\cite{shysheya2024conditional} revealed that this approximation hinders the propagation of local information beyond the temporal receptive field of the model. This limitation becomes problematic for systems governed by slowly evolving dynamics, where distant past observations strongly constrain present trajectories.\sidenote{\vspace{-4em}\textit{The perturbations which the motions of planets suffer from the influence of other planets, are so small and so slow that they only become sensible after a long interval of time.}\\[1ex] --- C. F. Gauss (1809) \nocite{gauss1809theoria}}

In forecasting tasks, where observations beyond the present-time cutoff are unavailable, \textcite{shysheya2024conditional} demonstrate that an autoregressive generation strategy effectively extends the range of information propagation. Their findings are consistent with our own forecasting experiments in Chapter \ref{chap:appa}. In the case of filtering, this approach is also applicable as past observations can be assimilated autoregressively. Conversely, autoregression is not immediately applicable to reanalysis, where past and future observations must be assimilated simultaneously.

Within the SDA framework, we foresee several approaches to address this persistent limitation. First, training models with wider temporal windows can improve assimilation quality, especially when observations are dense~\cite{shysheya2024conditional}. However, this strategy quickly becomes computationally expensive during training and inference. A related approach would be to coarsen the temporal resolution of the model by increasing the time stride between consecutive window elements. Such model would have to understand long-term dependencies to denoise effectively, while reducing the computational cost compared to the previous proposition. Unfortunately, with this approach, it becomes difficult to assimilate high-cadence observations, whose corresponding states are not present in a coarse trajectory. A model with heterogeneous windows, containing a mixture of dense-local and sparse-global elements, may strike the right balance between fine and coarse temporal resolutions, offering an efficient way to propagate local information globally. In practice, there is a risk that such model learns to ignore the global elements when denoising the local ones, defeating the original purpose.

Beyond architectural modifications, the path to accurate long-term modeling could also stem from algorithmic improvements. A recent line of work \cite{ge2023preserve, chang2024how, daras2024warped} demonstrates that using time-correlated noise, instead of white noise, can transform a pre-trained image diffusion model into a naive video model. Adapting this strategy to the SDA framework could mitigate the issue of long-term decorrelation, but is unlikely to solve it. From another perspective, \textcite{du2023reduce} describe several methods to compose the outputs of multiple diffusion models during generation. These methods are different from the one presented in Chapter \ref{chap:sda} and could lead to better information propagation. Finally, efficient global models could replace local models when the latter become too imprecise. Notably, in low signal-to-noise regimes, analytical denoisers derived from the data's mean and covariance~\cite{sarkka2013bayesian} can outperform local neural network-based denoisers, as they attend to the entire trajectory rather than finite windows. Implementing such an analytical denoiser is technically straightforward and should therefore be prioritized for investigation.

Together, these directions highlight a broader open question~\cite{beltagy2020longformer, tay2021long, gu2022efficiently, dwivedi2022long, su2024roformer} in deep learning: how to encourage long-range understanding in neural networks?

\vspace{-1ex}
\subsection*{Posterior approximations can be unfaithful}

In Bayesian inference, the posterior distribution encapsulates all the knowledge we possess about the problem at hand. If a posterior approximation is too inaccurate, it can lead to misleading conclusions that can propagate through downstream analysis and eventually influence policy or operational decisions. Therefore, posterior validation is not only a matter of scientific rigor, but a necessity to establish trust~\cite{hermans2022trust} in deep probabilistic models. Unfortunately, this aspect is often overlooked within the literature, and this dissertation is no exception.

The main reason for this frequent shortcoming is the overall lack of simultaneously principled and scalable diagnostics. Widespread metrics in the image generation literature, such as PSNR, SSIM~\cite{wang2004image}, and LPIPS~\cite{zhang2018unreasonable} are cheap to evaluate but poorly aligned with posterior faithfulness: PSNR and SSIM favor smooth reconstructions, whereas LPIPS emphasizes perceptual similarity. FID~\cite{heusel2017gans} is sound but applies only to unconditional generation, as the true conditional/posterior distribution is typically unavailable. Conversely, common posterior diagnostics in the Bayesian inference literature, such as expected coverage probability (ECP) of highest posterior density regions~\cite{rozet2021arbitrary, hermans2022trust, miller2021truncated, deistler2022truncated}, are rigorous but expensive in high dimensions, rely on probability density estimates, and fail to detect uninformative posteriors~\cite{hermans2022trust, lemos2023samplingbased}. Recent sample-based alternatives such as TARP~\cite{lemos2023samplingbased} or PQMass~\cite{lemos2025pqmass} appear to be more robust, but remain computationally demanding. In meteorology, the continuous ranked probability score (CRPS)~\cite{hersbach2000decomposition, gneiting2007strictly} is both principled and scalable, but only probes one-dimensional marginals. Multivariate scoring rules exist~\cite{gneiting2007strictly}, but their discriminative power is notoriously limited~\cite{pinson2013discrimination}. Finally, most metrics are meaningful only in expectation over multiple observations and, thus, not appropriate for case-by-case decision making. The local classifier two-sample test ($\ell$-C2ST) proposed by \textcite{linhart2023lc2st} enables localized posterior validation for individual observations, but relies on an auxiliary classifier trained with held-out calibration data and remains untested in high-dimensional regimes.

In other words, developing evaluation diagnostics that are accessible, principled, scalable, and useful remains an open challenge in making deep probabilistic models scientifically trustworthy and actionable.

\vspace{-1ex}
\subsection*{Mining gold from observations}

Beyond methodological and experimental shortcomings, an aspect missing from this dissertation is the integration of the DiEM algorithm, presented in Chapter \ref{chap:diem}, within the SDA framework, introduced in Chapter \ref{chap:sda}.\sidenote{\textit{Tant que l'Algèbre et la Géométrie ont été séparées, leurs progrès ont été lents et leurs usages bornés; mais lorsque ces deux sciences se sont réunies, elles se sont prêté des forces mutuelles et ont marché ensemble d'un pas rapide vers la perfection.}\\[1ex] --- J.-L. Lagrange (1795)} Combining the two could, in principle, enable learning global atmospheric models directly from observational data, thereby circumventing the need for synthetic data and hand-crafted models that typically misrepresent sub-resolution phenomena. This approach recalls the iterative strategy proposed by \textcite{bocquet2020bayesian}and \textcite{brajard2020combining, brajard2021combining}; where data assimilation acts as the expectation step, while training an emulator of the dynamics corresponds to the maximization step. In our work, we either rely on reanalysis data from ERA5~\cite{hersbach2020era5} or synthetic data from a simulator to train models. In both cases, our models inherit the biases of the underlying physical models, and become misspecified with respect to the true dynamics. Detecting and mitigating model misspecification in Bayesian inference is an active area of research \cite{cherief-abdellatif2020mmdbayes, wehenkel2022robust, dellaporta2022robust, cannon2022investigating, ward2022robust, huang2023learning, frazier2023reliable, gao2023generalized, nott2024bayesian, schmitt2024detecting, kelly2024misspecificationrobust, anaumontel2025tests, pierre2025mitigating, wehenkel2025addressing, senouf2025inductive, behrmann2025inferring} and refining misspecified diffusion models with DiEM is a promising approach, adopted by \textcite{barco2025tackling}.

Another long-standing obstacle in data assimilation, inherited by SDA, is to exploit observations for which no explicit and differentiable model of the measurement process exists, such as satellite data. Previous studies~\cite{andrychowicz2023deep, alexe2024graphdop, allen2025endtoend} have tackled this issue by training end-to-end models to predict one modality of observations from another modality. These purely data-driven approaches automatically compensate for unknown phenomena or biases in the data, at the expense of flexibility: incorporating a new type of observation typically requires retraining. This stands in contrast to the modular, plug-and-play nature of training-free guidance methods~\cite{chung2023diffusion, song2023pseudoinverseguided, boys2023tweedie}, central to the SDA framework.

An intermediate solution, compatible with SDA, would be to train differentiable surrogate models for the unknown measurement processes \cite{chevallier2000use} or their inverse. Such surrogates could then be naturally integrated into the perturbed likelihood approximations of Chapters \ref{chap:sda} or \ref{chap:appa}. However, training them requires paired state-observation data, which is rarely accessible in earth sciences. A possible remedy is to integrate this training step within the EM algorithm, using samples from the expectation step to refine surrogate measurement models, in the spirit of Gibbs sampling \cite{geman1984stochastic, heurtel-depeiges2024listening}. This \emph{blind} EM scheme could, in principle, jointly learn a prior model for the dynamics and surrogate models for the unknown measurement processes.

Whether through measurement surrogates, end-to-end models, or EM-like methods, considerable investments remain to be made for integrating a larger fraction of observational data within NWP pipelines.

\vspace{-1ex}
\section{Conclusion}

Ultimately, this thesis argues for a more probabilistic view of dynamical systems and scientific modeling; one that acknowledges uncertainty in observations, states, and models. This objective invites to move beyond the pitfalls of hand-crafted models and embrace data-driven methods, that can learn automatically without inheriting our biases. Our work takes a step forward in this direction by showing that generative models and, in particular, diffusion models can replicate, discover, and abstract the dynamics of our universe, solely from data. Many may oppose our position as they see machines as mere number jugglers, incapable of producing new knowledge. But, as Turing~\cite{turing1950computing} reminds us, there is value in the systematic derivation of consequences from data. Machines do so tirelessly, and the fruits of their labor surprise us with great frequency.

\clearpage
\section*{References}

\printbibliography[heading=none]
